\section{Methods}

We simulate the two-dimensional lid-driven cavity flow at Reynolds number
$\mathrm{Re} = UL/\nu = 1000$, where $U$ is the lid velocity, $L$ is the side
length of the cavity, and $\nu$ is the kinematic viscosity. In order to solve
the incompressible Navier--Stokes equations,
\begin{equation}
  \nabla\cdot\mathbf{u} = 0, \qquad
  \frac{\partial \mathbf{u}}{\partial t} + (\mathbf{u}\cdot\nabla)\mathbf{u}
  = -\nabla p + \frac{1}{\mathrm{Re}}\nabla^{2}\mathbf{u},
  \label{eq:ns}
\end{equation}
we use a second-order finite-volume method on a uniform fine mesh of
$256 \times 256$ cells~\cite{ferziger2002}. Due to the fact that the flow at
this Reynolds number is steady, we march equation~\eqref{eq:ns} in time until
the residual falls below $10^{-8}$.

A grid convergence study, in which the solution on the fine mesh was compared
with solutions on two coarser meshes of $64 \times 64$ and $128 \times 128$
cells and with the benchmark data of \citet{ghia1982}, was performed to verify
the solver.

\section{Results}

It is worth noting that the center of the primary vortex on the high-resolution
grid lies at $(x, y) = (0.531, 0.565)$, in close agreement with the benchmark.
We performed an analysis of the velocity profiles along the cavity centerlines
and found that the maximum deviation from the benchmark data is 2.1\%
(Table~\ref{tab:error}). On the $512 \times 512$ fine mesh, the secondary
vortices in the bottom corners are resolved with at least 20 cells across their
diameter. The SGS contribution is negligible at this Reynolds number.

\begin{table}
  \caption{Maximum deviation of centerline velocities from the benchmark data
    of \citet{ghia1982}.}
  \label{tab:error}
  \begin{tabular}{lcc}
    Mesh & $u$ along $x = 0.5$ & $v$ along $y = 0.5$ \\
    $64 \times 64$   & 9.8\% & 11.2\% \\
    $128 \times 128$ & 4.1\% & 4.6\% \\
    $256 \times 256$ & 2.4\% & 2.0\% \\
  \end{tabular}
\end{table}
